\chapter*{Declarations}

\paragraph{Use of Generative AI} I acknowledge the use of Gemini 3 (Google, \url{https://gemini.google.com/app}) to assist with writing Linux and LaTeX commands. It was used to help translate equivalence tests from Scala and C into Dafny. It was used to write a Dockerfile for my CI pipeline. It was used to provide suggestions to improve this report (though the report was entirely written by myself). I confirm that no content generated by AI has been presented as my own work.  

\paragraph{Ethical Considerations} EquiDafny must obey the licenses of any third-party libraries or tools it uses. In particular, our tool must obey Dafny's MIT license, Stainless's Apache 2.0 license\cite{StainlessFrontendsBenchmarks} and EqBench's MIT license\cite{badihiEqBenchDatasetEquivalent2021}. EquiDafny has a responsibility to be sound (it must not state that two programs are equivalent when they are not). EquiDafny is designed to be sound; in particular, the result of whether two programs are equivalent is decided by the Dafny program verifier, which is also designed to be sound~\parencite{FAQ}. However, this has not been formally proven.

\paragraph{Sustainability} EquiDafny was implemented and run entirely locally on a Lenovo IdeaPad with an Intel Core i7. However, as described in \Cref{chp:evaluation}, Dafny verification can consume significant amounts of CPU time (though this can be mitigated by imposing a timeout). Despite this, overall, energy consumption is minimal. 

\paragraph{Availability of Data and Materials}

The evaluation testsuite is publicly available at \url{https://github.com/NVictor2004/dafny-equivalence-benchmarks}. EquiDafny's source code is internally available within Imperial College London's dedicated GitLab instance at \url{https://gitlab.doc.ic.ac.uk/nv323/proving-program-equivalence}. 